% 图论（综述）
% license CCBYSA3
% type Wiki

本文根据 CC-BY-SA 协议转载翻译自维基百科\href{https://en.wikipedia.org/wiki/Graph_theory}{相关文章}

\begin{figure}[ht]
\centering
\includegraphics[width=6cm]{./figures/099d0167fdbd2081.png}
\caption{} \label{fig_TUl_1}
\end{figure}
在数学和计算机科学中，图论是对图的研究。图是一种数学结构，用于刻画对象之间的成对关系。在这一语境下，一个图由顶点（亦称为节点或点）组成，顶点之间通过边（亦称为弧、链接或线）连接。图论中区分无向图与有向图：在无向图中，边对两个顶点的连接是对称的；而在有向图中，边对两个顶点的连接是不对称的。图是离散数学研究的主要对象之一。
\subsection{定义}
图论中的定义不尽相同。以下是一些更基本的定义图及其相关数学结构的方式。
\subsubsection{图}
\begin{figure}[ht]
\centering
\includegraphics[width=6cm]{./figures/0674f0c8b9e7f86a.png}
\caption{一个具有三个顶点和三条边的无向图。} \label{fig_TUl_2}
\end{figure}
在一个受限但非常常见的意义下，[1][2]图是一个有序对$G = (V, E)$其中：
\begin{itemize}
\item $V$：顶点的集合（亦称为节点或点）；
\item $E \subseteq {{x, y} \mid x, y \in V ;\text{且}; x \neq y}$：边的集合（亦称为链接或线），其元素是无序的顶点对（即一条边关联两个不同的顶点）。
\end{itemize}
为了避免歧义，这类对象可称为无向单图。

在边${x, y}$中，顶点$x$与$y$称为该边的端点。该边被认为连接$x$与$y$，并与$x$、$y$关联。图中可能存在某些顶点没有属于任何边。根据该定义，不允许出现多重边（即两条或以上的边连接相同的一对顶点）。
\begin{figure}[ht]
\centering
\includegraphics[width=6cm]{./figures/e0f64ecda97f8953.png}
\caption{一个具有3个顶点、3条边和4条自环的无向多重图示例。} \label{fig_TUl_3}
\end{figure}
在另一种更一般的意义下，允许多重边，[3][4]图是一个有序三元组$G = (V, E, \phi)$其中：
\begin{itemize}
\item $V$：顶点的集合（也称为节点或点）；
\item $E$：边的集合（也称为链接或线）；
\item $\phi : E \to {{x,y} \mid x,y \in V ;\text{且}; x \neq y}$：一个关联函数，它将每一条边映射到一个无序的顶点对（即一条边连接两个不同的顶点）。
\end{itemize}
为了避免歧义，这类对象可称为无向多重图。

自环是一条将顶点连接到其自身的边。在前述两种定义下，图都不能有自环。因为一条将顶点$x$与自身连接的自环对应的边为${x,x} = {x}$，而它并不属于${{x,y}\mid x,y \in V, x \neq y}$。若要允许自环，则必须扩展定义：对于无向单图，其边集应修改为
$E \subseteq {{x,y}\mid x,y \in V}$；对于无向多重图，其关联函数应修改为
$\phi : E \to {{x,y}\mid x,y \in V}$。为了避免歧义，这些类型的对象可分别称为允许自环的无向单图和允许自环的无向多重图（有时也称为无向伪图 pseudograph）。

通常情况下，$V$和$E$被认为是有限的；许多著名的结论在无限图中并不成立（或有所不同），因为许多推理在无限情形下失效。此外，通常假设 (V) 是非空的，但 (E) 可以是空集。图的阶是$|V|$，即其顶点数。图的大小是$|E|$，即其边数。一个顶点的度（degree 或 valency）是与其关联的边的条数，其中自环计作两次。图的度是所有顶点度数的最大值。

在一个阶为$n$的无向单图中，每个顶点的最大度为$n-1$，该图的最大边数为$\frac{n(n-1)}{2}$.

允许自环的无向单图$G$的边集，在其顶点集上诱导出一个对称齐次关系$\sim$，称为$G$的邻接关系。具体而言，对于每条边$(x, y)$，其端点$x$和$y$被称为邻接，记作$x \sim y$.
\subsubsection{有向图}
\begin{figure}[ht]
\centering
\includegraphics[width=6cm]{./figures/95d6ef80502e422c.png}
\caption{一个具有三个顶点和四条有向边的有向图（双箭头表示在两个方向上各有一条边）。} \label{fig_TUl_4}
\end{figure}
有向图（directed graph 或 digraph）是一种边具有方向性的图。在一个受限但非常常见的意义下，[5] 有向图是一个有序对$G = (V, E)$其中：
\begin{itemize}
\item $V$：顶点的集合（也称为节点或点）；
\item $E \subseteq {(x,y) \mid (x,y) \in V^{2} ;\text{且}; x \neq y}$：边的集合（也称为有向边、定向链接、定向线、箭头或弧），其元素是顶点的有序对（即一条边关联两个不同的顶点，且方向不可交换）。
\end{itemize}
为了避免歧义，这类对象可称为有向单图。在集合论和图论中，$V^{n}$表示由集合$V$中元素组成的$n$元组的集合，即长度为$n$的有序序列，序列中的元素不必两两不同。

在一条从$x$指向$y$的边$(x,y)$中：顶点$x$和$y$被称为该边的端点；$x$称为该边的尾，$y$称为该边的头；这条边被认为连接$x$与$y$，并且与$x$、$y$关联。图中可能存在一些顶点并不属于任何边。边$(y,x)$称为边$(x,y)$的反向边。在上述定义下，不允许出现多重边，即尾和头都相同的两条或更多的边。

在另一种更一般的意义下，允许多重边，[5]有向图是一个有序三元组$G = (V, E, \phi)$其中：
\begin{itemize}
\item $V$：顶点的集合（也称为节点或点）；
\item $E$：边的集合（也称为有向边、定向链接、定向线、箭头或弧）；
\item $\phi : E \to {(x,y) \mid (x,y) \in V^{2} ;\text{且}; x \neq y}$：一个关联函数，它将每条边映射到一个顶点的有序对（即一条边连接两个不同的顶点，并具有方向）。
\end{itemize}
为了避免歧义，这类对象可称为有向多重图。

自环是一条将顶点连接到其自身的边。在前述两种定义下的有向图中，自环是不允许的。因为一条将顶点 (x) 与自身相连的自环对应的边为$(x,x)$，但$(x,x)$并不属于${(x,y)\mid (x,y)\in V^{2}, ; x \neq y}$.因此，为了允许自环，必须扩展定义：对于有向单图，其边集的定义应修改为$E \subseteq {(x,y)\mid (x,y)\in V^{2}}$;对于有向多重图，其关联函数的定义应修改为$\phi : E \to {(x,y)\mid (x,y)\in V^{2}}$.为了避免歧义，这些对象可分别称为：允许自环的有向单图和允许自环的有向多重图（或称quiver，箭图）。

允许自环的有向单$G$的边集在其顶点集上诱导出一个齐次关系$\sim$，称为$G$的邻接关系。具体而言，对于每条边$(x,y)$，其端点$x$和$y$被称为邻接，记作$x \sim y$.
\subsection{应用}
\begin{figure}[ht]
\centering
\includegraphics[width=6cm]{./figures/5b073a7da7b89602.png}
\caption{由维基百科编辑者（边）与不同语言版本的维基百科（顶点）在2013年夏季某一个月内的贡献所形成的网络图。[6]} \label{fig_TUl_5}
\end{figure}
图可以用来建模物理、生物学[7][8]、社会和信息系统[9]中的多种关系与过程。许多实际问题都能用图来表示。为了强调它们在现实系统中的应用，网络这一术语有时被定义为：在图的顶点和边上附加属性（例如名称）的图。而把现实世界系统表达并理解为网络的学科被称为网络科学。
\subsubsection{计算机科学}
在计算机科学中，“因果”与“非因果”的链接结构就是用来表示通信网络、数据组织、计算设备、计算流程等的图。例如，一个网站的链接结构可以用有向图来表示，其中顶点（节点）代表网页，有向边代表从一个页面到另一个页面的超链接。类似的方法也可应用于社交媒体[10]、交通出行、生物学、芯片设计、神经退行性疾病进展的映射[11][12]以及许多其他领域。因此，开发处理图的算法是计算机科学中的一个重要研究方向。图的变换常常被形式化为图重写系统。与专注于基于规则的图内存操作的图变换系统相互补充的是图数据库，它们致力于对图结构数据进行事务安全、持久化的存储与查询。
\subsubsection{语言学}
图论方法以多种形式被证明在语言学中特别有用，因为自然语言往往很适合用离散结构来刻画。传统上，句法和组合语义学遵循基于树的结构，其表达能力依赖于组合性原则，并以层级化的图来建模。较为当代的方法，例如以词为中心的短语结构语法，则使用类型特征结构来建模自然语言的句法，而这类结构本质上是有向无环图。在词汇语义学中，尤其是在计算机应用中，若一个词的意义通过与相关词的关系来理解，其建模会更加容易；因此，语义网络在计算语言学中占有重要地位。除此之外，在音系学中（例如使用格图的最优化理论）以及形态学中（例如使用有限状态转换器的有限状态形态学），将语言分析为图的方式也十分常见。实际上，这一数学领域对语言学的价值催生了诸如TextGraphs这样的组织，以及一系列 “Net” 项目，例如WordNet、VerbNet等。
\subsubsection{物理学与化学}
图论也被用于研究化学和物理中的分子。在凝聚态物理中，复杂的三维原子结构模拟可以通过收集与原子拓扑相关的图论性质的统计量来进行定量研究。与此同时，“费曼图及其计算规则以一种与实验数据密切相关的形式总结了量子场论”[13]。在化学中，图为分子提供了一种自然的建模方式，其中顶点代表原子，边代表化学键。这种方法特别适用于分子结构的计算机处理，从化学编辑器到数据库搜索均有应用。在统计物理中，图可以表示系统中相互作用部分之间的局部连接，以及物理过程在这些系统中的动力学演化。类似地，在计算神经科学中，图可以用来表示大脑区域之间的功能连接，这些区域的相互作用产生各种认知过程，其中顶点表示不同的大脑区域，边表示这些区域之间的连接。图论在电学网络的建模中也扮演着重要角色：此时边带有权重，该权重与导线段的电阻相关，用于获得网络结构的电学性质[14]。此外，图还可用于表示多孔介质的微观通道，其中顶点代表孔隙，边代表连接孔隙的细小通道。化学图论使用分子图来对分子建模。图与网络是研究和理解相变与临界现象的极佳模型。移除节点或边会导致一个临界转变，使得网络破裂成更小的簇，这被视为一种相变。该现象通过渗流理论进行研究[15]。
\subsubsection{社会科学}
\begin{figure}[ht]
\centering
\includegraphics[width=6cm]{./figures/0c45309aaeff9f1e.png}
\caption{} \label{fig_TUl_6}
\end{figure}
图论在社会学中也被广泛应用，例如用来衡量行动者的声望，或研究谣言传播，尤其是通过社会网络分析软件来实现。在社会网络这一范畴下，存在多种不同类型的图。[17]熟人图与朋友图：描述人们之间是否互相认识。影响图：建模某些人是否能够影响他人的行为。合作图：建模两个人是否以某种特定方式进行合作，例如共同出演一部电影。
\subsubsection{生物学}
同样，图论在生物学和保护研究中也非常有用。在这种应用中，**顶点**可以表示某些物种存在（或栖息）的区域，**边**则表示区域之间的迁徙路径或移动。这类信息在研究繁殖模式，或追踪疾病、寄生虫传播，以及考察迁徙变化如何影响其他物种时，具有重要意义。

图论还常用于**分子生物学**和**基因组学**，以建模和分析具有复杂关系的数据集。例如，在**单细胞转录组分析**中，基于图的方法常被用于将细胞“聚类”为不同的细胞类型。另一种应用是对通路中的基因或蛋白质建模，并研究它们之间的关系，如**代谢通路**和**基因调控网络**[18]。此外，**进化树**、**生态网络**以及**基因表达模式的层次聚类**也常用图结构来表示。

图论还应用于**连接组学（connectomics）**[19]：神经系统可以看作一张图，其中**节点**是神经元，**边**则是它们之间的连接。
\subsubsection{数学}
在数学中，图在**几何学**和**拓扑学的某些部分**（如结理论）中非常有用。**代数图论**与群论关系密切。代数图论已被应用于许多领域，包括**动力系统**和**复杂性研究**。
\subsubsection{其他主题}
图结构可以通过为每条边赋予一个权重来扩展。带权图（weighted graphs）用于表示那些成对连接带有数值信息的结构。例如，如果图表示一个道路网络，边的权重可以表示每条道路的长度。每条边也可能带有多个权重，如距离（如前例所示）、行驶时间或金钱成本。这类带权图常用于**GPS 程序**和**旅行规划搜索引擎**，例如比较航班时间和费用。
